% Options for packages loaded elsewhere
\PassOptionsToPackage{unicode}{hyperref}
\PassOptionsToPackage{hyphens}{url}
%
\documentclass[
]{article}
\usepackage{amsmath,amssymb}
\usepackage{iftex}
\ifPDFTeX
  \usepackage[T1]{fontenc}
  \usepackage[utf8]{inputenc}
  \usepackage{textcomp} % provide euro and other symbols
\else % if luatex or xetex
  \usepackage{unicode-math} % this also loads fontspec
  \defaultfontfeatures{Scale=MatchLowercase}
  \defaultfontfeatures[\rmfamily]{Ligatures=TeX,Scale=1}
\fi
\usepackage{lmodern}
\ifPDFTeX\else
  % xetex/luatex font selection
\fi
% Use upquote if available, for straight quotes in verbatim environments
\IfFileExists{upquote.sty}{\usepackage{upquote}}{}
\IfFileExists{microtype.sty}{% use microtype if available
  \usepackage[]{microtype}
  \UseMicrotypeSet[protrusion]{basicmath} % disable protrusion for tt fonts
}{}
\makeatletter
\@ifundefined{KOMAClassName}{% if non-KOMA class
  \IfFileExists{parskip.sty}{%
    \usepackage{parskip}
  }{% else
    \setlength{\parindent}{0pt}
    \setlength{\parskip}{6pt plus 2pt minus 1pt}}
}{% if KOMA class
  \KOMAoptions{parskip=half}}
\makeatother
\usepackage{xcolor}
\setlength{\emergencystretch}{3em} % prevent overfull lines
\providecommand{\tightlist}{%
  \setlength{\itemsep}{0pt}\setlength{\parskip}{0pt}}
\setcounter{secnumdepth}{-\maxdimen} % remove section numbering
\ifLuaTeX
  \usepackage{selnolig}  % disable illegal ligatures
\fi
\usepackage{bookmark}
\IfFileExists{xurl.sty}{\usepackage{xurl}}{} % add URL line breaks if available
\urlstyle{same}
\hypersetup{
  pdftitle={Fixed-root Higgs shifts in the non-Abelian five-point tree},
  hidelinks,
  pdfcreator={LaTeX via pandoc}}

\title{Fixed-root Higgs shifts in the non-Abelian five-point tree}
\author{}
\date{}

\begin{document}
\maketitle

\section{\texorpdfstring{When do moving poles suffice? The non-Abelian
(W\textsuperscript{+W}-W\^{}3hh)
tree}{When do moving poles suffice? The non-Abelian (W+W-W\^{}3hh) tree}}\label{when-do-moving-poles-suffice-the-non-abelian-ww-w3hh-tree}

\subsection{Result and conventions}\label{result-and-conventions}

For three physical equal-mass vectors in a renormalizable single-doublet
(SU(2)) theory, the complex two-Higgs shift gives a bounded, generically
nonzero tree amplitude. The exact boundary separates into (i) the
seagull and stationary Higgs exchange, (ii) ordered emissions on a
single vector arm, and (iii) emissions from different arms of a true
cubic gauge vertex. The first two use the same fixed-root building
blocks as the
\href{../su2-fixed-root-boundary/OffShell_Polarization_Boundary.md}{neutral
result}. The third cannot be obtained by sewing just one neutral
two-source boundary. We give a computable zero-boundary criterion and a
nonzero stationary-pole witness. This is a result for the specified
shift and field content, not a statement about all electroweak
amplitudes.

All momenta are incoming, metric (+---); (k\_i\textsuperscript{2=m}2),
(e\_i\cdot k\_i=0), (p\textsuperscript{2=r}2=M\^{}2), and
(k\_1+k\_2+k\_3+p+r=0). Here (m) is the common vector mass, (M=m\_h),
(a=g\_2/2), and (v=m/a). Work in unitary gauge and factor out a common
cubic-group factor and overall phase. In the real adjoint basis one may
use legs (1,2,3); changing to (+,-,3) only changes the common color
phase, because each surviving tree contains exactly one antisymmetric
cubic group tensor, and scalar vertices and propagators are color
diagonal. The cubic kinematic vertex is {[}
V(k\_1,k\_2,k\_3;e\_1,e\_2,e\_3)
=\sum\emph{\{\rm cyc\}(e\_i\cdot e\_j)(k\_i-k\_j)\cdot e\_s. {]} Use
(U(Q)e={[}e-Q(Q\cdot e)/m\^{}2{]}/D(Q)), (D(Q)=Q\textsuperscript{2-m}2),
(g\_h=2am), (g}\{hh\}=2a\^{}2). Let (H=p+r) and shift
(p\mapsto p+zn,~r\mapsto r-zn), with (n\^{}2=n\cdot p=n\cdot r=0).
Require (\ell\_i=2k\_i\cdot n\ne0) and nondegenerate propagators;
complex contractions are bilinear with no conjugation. The space of
projective null directions has two choices at generic equal-mass (p,r)
kinematics.

\subsection{Complete graph count and full
amplitude}\label{complete-graph-count-and-full-amplitude}

A five-leaf tree has either three cubic vertices or one cubic and one
quartic. The only vertices of unitary-gauge single-doublet (SU(2)) are
(WWW), (WWWW), (WWh), (WWhh), (hhh), (hhhh). An odd number of external
vectors requires a cubic (WWW). There are \textbf{18 labelled trees}:

\begin{itemize}
\tightlist
\item
  6 ordered emissions of both Higgs particles on a single arm;
\item
  6 placements on different arms;
\item
  3 (WWhh) seagulls;
\item
  3 (h\^{}*\to hh) splits.
\end{itemize}

The quartic gauge vertex cannot attach to two Higgs particles through
one allowed cubic vertex. A separate partition recursion with all these
vertex types and fixed external adjoint colors finds the same count.

For each arm (i), define (K\_i=k\_i+H) and {[}
C\_0=2a\^{}2\left(1+\frac{3M^2}{H^2-M^2}\right),\qquad
E\_i(p,r)=U(K\_i)\{C\_0e\_i+g\_h\^{}2{[}U(k\_i+p)+U(k\_i+r){]}e\_i\}.
{]} A term with both Higgs particles on arm (i) is (V) with the (i)-th
arguments replaced by (K\_i,E\_i). For distinct arms (i,j), replace
those two arguments by (k\_i+x,g\_hU(k\_i+x)e\_i) and
(k\_j+y,g\_hU(k\_j+y)e\_j), summing over cyclic arm pairs and
((x,y)=(p,r),(r,p)). This is the complete signed amplitude before taking
any large-(z) limit.

\subsection{Closed boundary, and a practical cancellation
test}\label{closed-boundary-and-a-practical-cancellation-test}

Set (R\_i=(n\cdot e\_i)/\ell\emph{i) and {[}
F\_i=\frac{H\cdot e_i-2R_i(M^2+k_i\cdot H)}{\ell_i},\qquad
S\_i=-\frac{R_i(2k_i+H)+nF_i}{m^2}. {]} For the same-arm terms, expand
(U(k\_i+x+wn)e\_i) for (w=\pm z). Its linear part is (-w,nR\_i/m\^{}2),
which cancels between the two orders. The sum of constants is (S\_i).
Thus the same-arm plus fixed part (B}\{\rm same\}) is obtained by
inserting {[} B\_i=U(K\_i)(C\_0e\_i+g\_h\^{}2 S\_i) {]} in the (i)-th
arm of (V), again using momentum (K\_i). The other two physical
equal-mass vectors yield a conserved (WWW) current at that arm, as
follows directly from contracting the cubic vertex with its root
momentum. This is why the neutral fixed-root two-source construction can
be sewn into this part.

For different arms, let (L\_i=k\_i+x+wn,~L\_j=k\_j+y-wn),
(L\_i+L\_j+k\_s=0), (b\_i=L\_i\cdot e\_i,~b\_j=L\_j\cdot e\_j). The two
contractions with (L\_i) and (L\_j), using (k\_s\textsuperscript{2=m}2)
and (k\_s\cdot e\_s=0), yield the exact identity {[}

\begin{split}
 V(L_i,L_j,k_s;U(L_i)e_i,U(L_j)e_j,e_s)
 &= \frac{V(L_i,L_j,k_s;e_i,e_j,e_s)}{D_iD_j}\\
 &+\frac{1}{m^2}\left[\frac{b_i}{D_i}e_j\cdot e_s
   -\frac{b_j}{D_j}e_i\cdot e_s
   +\frac{b_ib_j}{D_iD_j}L_i\cdot e_s\right].
\end{split}

{]} Here (D\_\{i,j\}=D(L\_\{i,j\})). The first term decays as (1/w); the
last term seems to grow as (w), but its growth cancels when the labelled
Higgs assignments are paired. The two single-ratio constants
proportional to (R\_i(e\_j\cdot e\_s)) and (R\_j(e\_i\cdot e\_s)) cancel
in the cyclic sum. The surviving cross-arm boundary is {[}
B\_\{\rm cross\}=\frac{g_h^2}{m^2}\sum\emph{\{\rm cyc\}
\left[ R_iR_j(k_i-k_j)\cdot e_s+
 (n\cdot e_s)(R_jF_i-R_iF_j)\right]. {]} There is no growing part of the
sum. The boundary is exactly
(B}\infty=B\_\{\rm same\}+B\_\{\rm cross\}). Given ((k\_i,e\_i,p,r,n)),
these formulas use the unshifted vectors and scalar invariants and no
large-(z) fit. The moving physical poles alone recover (A(0)) \textbf{if
and only if} this sum vanishes at the point. In particular:

\begin{itemize}
\tightlist
\item
  If at least two physical polarizations obey (n\cdot e\_i=0), then
  (B\_\{\rm cross\}=0) identically. One such polarization does not
  suffice. The stationary and same-arm terms need not vanish.
\item
  The \emph{full} vanishing condition is a linear kernel in any one
  vector polarization when the other two are fixed. The explicit
  controls include nonzero amplitudes lying in this kernel, rather than
  setting the entire process to zero. The kernel is over complex
  physical polarizations; it is not in general a statement about a
  standard helicity eigenstate.
\end{itemize}

The boundary contains a physical stationary pole in (H\^{}2), so the
full boundary cannot vanish identically for either regular branch of the
shift. Its residue at (H\textsuperscript{2=M}2) is ((3M\^{}2/v)
A\_4(W\textsuperscript{+W}-W\^{}3h)). The other terms are generically
regular at that pole and cannot cancel it. For a concrete witness,
choose opposite incoming three-momenta along (z), outgoing vector and
scalar along (x), the first two polarizations along (y) and the third
along (z). Then (A\_4=2g\_h\textbar{}\mathbf k\_1\textbar{}\sum\_i
1/D(K\_i)); this is not identically zero, as seen also by its
large-energy coefficient (-3/(4 E\^{}2)) in the denominators' sum. The
physical five-point (hh) region has (H\textsuperscript{2\ge4M}2); the
pole argument uses analytic continuation, not a claimed real (h\to hh)
decay.

\subsection{Independent checks}\label{independent-checks}

The recorded 40 deformations (20 real (2\to3) phase-space points, two
null directions) have six distinct moving vector poles. For each pole
the three analytically dual spin states sew an (WWh) vertex to the
complete four-point (WWWh) tree, and agree with a small-contour
extraction to (8.9\times10\^{}\{-14\}) scaled. The independent rooted
tree evaluator uses explicit color permutations of the derivative slots,
rather than the grouped formula; agreement is (2.6\times10\^{}\{-15\}).
A contour outside all moving poles checks the constant to
(2.2\times10\^{}\{-13\}), and restoring it reconstructs the unshifted
result to (7.9\times10\^{}\{-15\}). The output also includes
positive-power contour checks, two-polarization cancellation controls
that retain a nonzero full boundary, and full-boundary cancellation
controls that retain a nonzero amplitude. No finite-(\xi) tensor
cross-check is claimed for the charged process.

\subsection{What is and is not new relative to known
results}\label{what-is-and-is-not-new-relative-to-known-results}

\begin{itemize}
\tightlist
\item
  \href{https://arxiv.org/html/2611.13096v1\#S3}{Ke--Liu--Mahbub,
  arXiv:2611.13096v1 §§3--4, Appendix B} derive a polarization
  continuation from on-shell factorization and use ALT for their
  constructibility result. Our fixed-root scalar shift leaves some
  propagator invariants constant and gives \emph{linear} moving
  denominators; it does not obey the large-(z) assumptions of their ALT
  bound. This is not an ALT counterexample or a different physical
  massive spin sum.
\item
  \href{https://arxiv.org/html/2403.15538\#S2.SS2}{Ema et al.,
  arXiv:2403.15538, §2.2} explicitly distinguish their ALT from null
  shifts whose propagators scale differently, and discuss spin-dependent
  constructibility. Shift dependence itself is not a novelty here. We
  supply explicit gauge/Higgs boundary coefficients and polarization
  cancellation conditions for one different five-point deformation.
\item
  \href{https://arxiv.org/abs/1910.13407}{Franken--Schwinn,
  arXiv:1910.13407, §5} already analyze massive tree constructibility
  and when more lines suffice. Their existence and scaling statements do
  not identify the fixed-vector two-Higgs boundary written above.
\item
  \href{https://doi.org/10.1007/JHEP12\%282019\%29058}{Durieux,
  Kitahara, Shadmi and Weiss, JHEP 12 (2019) 058, §5.2} explicitly
  discuss adverse large-(z) behavior and reconstruct the four-point
  (W\overline W hh) amplitude from all physical invariants plus a
  nonfactorizable term determined by the ultraviolet behavior. This is
  close precedent: neither the existence of a Higgs contact nor failure
  of a two-line shift at four points is our novelty. The present result
  adds the two moving legs of the cubic gauge vertex, their cross-arm
  term and its polarization selection rule, and distinguishes poles
  stationary in (z) from interactions local in all invariants at five
  points.
\item
  \href{https://arxiv.org/abs/1507.00463}{Jin--Feng, arXiv:1507.00463}
  and \href{https://arxiv.org/abs/2210.07025}{Huang--Jin--Li,
  arXiv:2210.07025} supply systematic boundary-operator methods. We do
  not claim to invent boundary operators or the rational contour
  identity. Our scoped contribution is the explicit decomposition and
  the algebraic two-polarization selection rule for a genuine
  broken-gauge five-point process, paired with the neutral two-source
  and Yukawa six-point checks.
\end{itemize}

The case is no longer merely one more example of a nonzero boundary: it
resolves what must be added, what vanishes as a tensor contraction, and
why even the simpler common contribution can survive. It is still a
narrow result. We do not claim general-multiplicity, nondegenerate
electroweak, complex-pole or loop theorems, or guaranteed publication.

\end{document}
